\documentclass[11pt,a4paper]{article}

\usepackage[utf8]{inputenc}
\usepackage[T1]{fontenc}
\usepackage{amsmath,amssymb,amsthm}
\usepackage{booktabs}
\usepackage{enumitem}
\usepackage{microtype}
\usepackage[margin=2.5cm]{geometry}
\usepackage{hyperref}
\hypersetup{colorlinks=true,linkcolor=blue,citecolor=blue}

\title{From KGQA to Constrained Generation:\\
A Positioning of DCA-Trie in the Knowledge Graph\\
Question Answering Landscape}
\author{Bernard Katamanso}
\date{\today}

\begin{document}
\maketitle

%======================================================================
\section{What Is KGQA?}
%======================================================================

Knowledge Graph Question Answering (KGQA) is the task of answering a
natural language question $q$ using a structured knowledge graph
$\G = (V, E, \mathcal{R})$, where $V$ is a set of entities, $E \subseteq
V \times \mathcal{R} \times V$ is a set of labeled edges, and
$\mathcal{R}$ is a set of relation types.

Formally, given a question $q$ and a knowledge graph $\G$, the goal is
to produce an answer $a \in V$ (or a set of answers) such that $a$ is
the correct response to $q$ according to the information in $\G$.

\subsection{Why It Matters}

KGQA sits at the intersection of three converging trends:

\begin{enumerate}
  \item \textbf{Structured knowledge.}  KGs like Freebase, Wikidata, and
    YAGO encode billions of facts as $(h, r, t)$ triples.  They are
    auditable, updatable, and free of the parametric knowledge drift
    that plagues LLMs.

  \item \textbf{Natural language interface.}  Users want to query
    structured data in natural language, not SPARQL.  KGQA bridges this
    gap.

  \item \textbf{Reasoning over facts.}  Unlike simple fact lookup, KGQA
    requires \emph{multi-hop reasoning}: following a path through the
    graph to arrive at an answer that is not directly stated in the
    question.
\end{enumerate}

%======================================================================
\section{The KGQA Taxonomy}
%======================================================================

KGQA methods fall into four paradigms.  Each makes different assumptions
about how to bridge language and graph structure.

\subsection{Semantic Parsing}

Convert the question into a formal query (SPARQL, logical form) and
execute it against the KG.

\[
  q \;\longrightarrow\; \text{SPARQL}(q) \;\longrightarrow\;
  \text{execute}(\text{SPARQL}(q), \G) \;\longrightarrow\; a
\]

\textbf{Strengths:} precise, interpretable, verifiable.\\
\textbf{Weaknesses:} brittle to paraphrase, requires annotated
logical forms, struggles with ambiguous questions.

\subsection{Information Retrieval}

Retrieve a subgraph $\G_q \subseteq \G$ relevant to $q$, then answer
from the subgraph.

\[
  q \;\longrightarrow\; \text{retrieve}(q, \G) = \G_q
  \;\longrightarrow\; \text{reason}(\G_q) \;\longrightarrow\; a
\]

\textbf{Strengths:} scalable, handles large KGs.\\
\textbf{Weaknesses:} retrieval quality is the bottleneck; errors in
$\G_q$ propagate to the answer.

\subsection{Embedding-Based}

Learn vector representations of entities and relations, then answer by
vector similarity or geometric reasoning.

\[
  q \;\longrightarrow\; \text{encode}(q) = \mathbf{q}
  \;\longrightarrow\; \arg\max_{a \in V}\;
  \text{sim}(\mathbf{q}, \mathbf{a}) \;\longrightarrow\; a
\]

\textbf{Strengths:} smooth generalization, handles unseen entities.\\
\textbf{Weaknesses:} opaque reasoning, requires large-scale training,
struggles with multi-hop paths.

\subsection{Generation-Based (LLM + KG)}

Use a large language model to generate a reasoning path, constrained by
KG structure.

\[
  q,\; \G \;\longrightarrow\;
  \text{LLM}(q, \G_q) = (v_0, e_1, v_1, \ldots, e_L, v_L)
  \;\longrightarrow\; v_L = a
\]

\textbf{Strengths:} interpretable paths, natural language output,
leveraging LLM world knowledge.\\
\textbf{Weaknesses:} hallucination risk, computationally expensive.

\paragraph{This is where DCA-Trie sits.}  We extend this paradigm by
constraining the LLM's generation at \emph{every token} using a trie
derived from the KG, augmented with semantic type gates.

%======================================================================
\section{The Hallucination Problem}
%======================================================================

The fundamental challenge in generation-based KGQA is
\textbf{hallucination}: the LLM generates plausible-sounding entities
or relations that do not exist in the KG.

Consider:
\begin{quote}
  \textbf{Q:} ``Who directed Inception?''\\
  \textbf{LLM (unconstrained):} ``Inception was directed by
  Christopher Nolan, who also made The Dark Knight.''
\end{quote}

The second clause is factually correct but not grounded in the KG
subgraph for this question.  In a constrained setting, the model should
only output paths that exist in $\G$.

\subsection{Levels of Constraint}

\begin{center}
\begin{tabular}{@{}lll@{}}
\toprule
  Level & Mechanism & What it prevents \\
\midrule
  Token-level  & Trie constrains next token     & Hallucinated tokens \\
  Topological  & Edges must exist in $\G$        & Hallucinated edges \\
  Semantic     & Type gates enforce ontology     & Type-incoherent paths \\
  Query-aware  & Answer types from question      & Irrelevant answers \\
\bottomrule
\end{tabular}
\end{center}

GCR (rmanluo et al.\ 2025) operates at levels 1--2.  DCA-Trie adds
levels 3--4.

%======================================================================
\section{Why Graph-Constrained Reasoning?}
%======================================================================

The key insight is that KGQA is \emph{not} an open-ended generation
task.  The answer must be a real entity connected to the question entity
via a chain of real edges.  This structural constraint can be exploited
at decode time.

\subsection{The GCR Framework}

Graph-Constrained Reasoning (GCR)~\cite{gcr2025} introduced the idea of
encoding all valid paths from the KG into a trie, then using
\texttt{prefix\_allowed\_tokens\_fn} to restrict the LLM's vocabulary
at every decode step.

\paragraph{Baseline (unconstrained).}  The LLM generates freely:
\[
  P(v_L \mid q) = \prod_{l=1}^{L} P(v_l \mid v_{<l}, q).
\]

\paragraph{GCR (constrained).}  At each step $l$, the valid next tokens
are restricted to those that appear in some path prefix in the trie:
\[
  P(v_L \mid q, \mathcal{T}) = \prod_{l=1}^{L}
  P(v_l \mid v_{<l}, q,\; v_l \in \mathcal{T}(v_{<l})).
\]
where $\mathcal{T}(p)$ returns the set of valid next tokens given prefix
$p$.

This ensures every generated path is a real path in $\G$.

\subsection{The DCA-Trie Extension}

DCA-Trie adds semantic constraints on top of the topological constraint.
Before inserting paths into the trie, we filter by type compatibility:

\[
  \mathcal{T}_{\text{filtered}} = \{(h, r, t) \in \mathcal{T} \mid
  \text{range\_gate}(r, t) \;\land\;
  \text{type\_gate}(t, A, l, L)\}
\]

This is a \emph{pre-processing} step that reduces the trie size without
changing the decoding algorithm.

%======================================================================
\section{Complexity Analysis}
%======================================================================

\subsection{The Path Explosion Problem}

For a KG with average degree $d$ and path length $L$, the number of
possible paths from a query entity is:
\[
  |\text{paths}| = O(d^L).
\]

For WebQSP ($d \approx 20$, $L = 2$): $d^L \approx 400$ paths per
entity.  With 1,628 questions, this is manageable.

For CWQ ($d \approx 20$, $L = 4$): $d^L \approx 160{,}000$ paths.
This is the path explosion problem.

\subsection{How Each Method Addresses It}

\begin{center}
\small
\begin{tabular}{@{}lcc@{}}
\toprule
  Method & Trie size & Per-hop cost \\
\midrule
  GCR baseline   & $O(d^L)$ static   & $O(1)$ per token \\
  GCR + filtering & $O(d^L)$ smaller  & $O(1)$ per token \\
  DCA v1 (static) & $O(d^L)$ filtered & $O(1)$ per token \\
  DCA v2 (dynamic) & $O(d)$ per hop   & $O(1)$ per token \\
  DoG             & $O(d)$ per hop    & $O(1)$ per token \\
\bottomrule
\end{tabular}
\end{center}

The key advantage of dynamic methods (v2, DoG) is that the trie size is
independent of $L$: at each hop, we only enumerate 1-hop neighbors,
keeping the trie to $O(d)$ entries.

%======================================================================
\section{Benchmarks and Evaluation}
%======================================================================

\subsection{Standard Benchmarks}

\begin{itemize}
  \item \textbf{WebQSP}~\cite{webqsp}: 5,810 questions, 2--3 hop paths,
    Freebase subgraphs.  Simpler, shorter paths.
  \item \textbf{ComplexWebQuestions (CWQ)}~\cite{cwq}: 27,000+ questions,
    2--5 hop paths, more complex compositions.
  \item \textbf{MetaQA}~\cite{metaqa}: Movies domain, 1--3 hop questions.
  \item \textbf{QALD-9}: Cross-lingual, SPARQL-based.
\end{itemize}

\subsection{Our Evaluation}

\begin{center}
\begin{tabular}{@{}lccc@{}}
\toprule
  Method & WebQSP ($N$=100) & CWQ ($N$=500) & Path reduction \\
\midrule
  GCR baseline      & 89.0\% & 53.2\% & --- \\
  DCA v1 (filtered) & 88.0\% & ---    & 13.3\% \\
  DCA v2 (dynamic)  & 54.9\% & 32.2\% & --- \\
\bottomrule
\end{tabular}
\end{center}

%======================================================================
\section{Positioning DCA-Trie}
%======================================================================

\subsection{Relation to Prior Work}

\begin{center}
\small
\begin{tabular}{@{}lll@{}}
\toprule
  Work & Constraint type & Semantic filtering \\
\midrule
  GCR (2025)           & Topological (trie)      & No \\
  DoG (2025)           & Topological (dynamic)   & No \\
  KBQA-o1 (2025)       & Search-based (MCTS)     & No \\
  ORT (2024)           & Label-level reasoning   & Ontology-level \\
  \textbf{DCA-Trie}    & \textbf{Topological + semantic} & \textbf{Yes} \\
\bottomrule
\end{tabular}
\end{center}

\subsection{Contributions}

\begin{enumerate}
  \item \textbf{TypeOracle.}  A purely symbolic semantic gate using
    Freebase ontology (entity types, relation ranges).  No embeddings,
    no learned parameters.

  \item \textbf{Static filtering (v1).}  Pre-filter the trie before
    decoding.  Simple, effective for short paths.

  \item \textbf{Dynamic filtering (v2).}  Per-hop trie expansion with
    TypeOracle gates, following DoG's iterative architecture.

  \item \textbf{Empirical analysis.}  Systematic evaluation showing
    that filtering reduces paths by 13\% at 1pp accuracy cost (WebQSP),
    but the cost-benefit tradeoff degrades on deeper graphs (CWQ).
\end{enumerate}

\subsection{Limitations and Future Directions}

\begin{itemize}
  \item The Freebase ontology is broad---most entities have multiple
    types, so type gates are permissive.  Finer-grained ontologies
    (e.g., domain-specific) would improve pruning.

  \item Reasoning context (D4) is underexploited.  The generated relation
    itself constrains the next entity; this could be used for on-the-fly
    filtering without waiting for the type oracle.

  \item Multi-hop questions ($L \geq 3$) remain challenging for all
    methods.  The compounding error problem in iterative decoding is
    not yet solved.
\end{itemize}

%======================================================================
% Bibliography
%======================================================================
\begin{thebibliography}{1}

\bibitem{gcr2025}
  Luo et al., ``Graph-Constrained Reasoning: Faithful Reasoning on
  Knowledge Graphs with Large Language Models,'' \emph{ICML}, 2025.

\bibitem{dog2025}
  Lee et al., ``Decoding on Graph: Augmenting LLM Reasoning with
  Knowledge Graph Structured Search,'' \emph{ACL}, 2025.
  arXiv:2410.18415.

\bibitem{webqsp}
  Berant et al., ``Semantic Parsing on Freebase from Question-Answer
  Pairs,'' \emph{EMNLP}, 2013.

\bibitem{cwq}
  Talmor and Berant, ``The Web as a Knowledge-Base for Answering
  Complex Questions,'' \emph{NAACL}, 2018.

\bibitem{metaqa}
  Liang et al., ``Variational Reasoning for Question Answering with
  Knowledge Graph,'' \emph{AAAI}, 2018.

\bibitem{kbqa_o1}
  ``KBQA-o1: Agentic Knowledge Base Question Answering with Monte Carlo
  Tree Search,'' \emph{ICML}, 2025.

\bibitem{ort}
  Liu et al., ``ORT: Ontology-Guided Reverse Thinking for Knowledge
  Graph Question Answering,'' \emph{ACL}, 2025.

\end{thebibliography}

\end{document}
